% Section 6 -- Evaluation
\section{Evaluation}
\label{sec:eval}

The evaluation asks three questions. First, does the Phase-1 arbitration
assign the right cause class, and is its margin over a single-signal
baseline significant (\S\ref{sec:eval-cause})? Second, what are the
operating characteristics of the Phase-2 direction test---its size, its
power against planted mechanisms, and its coverage on the caches in the
tree---and does it commit a direction only when the evidence supports one
(\S\ref{sec:eval-direction})? Third, how do the two phases behave when the
input changes shape, when the event population changes, and on an
uncontrolled live alert stream (\S\ref{sec:eval-shape}--\S\ref{sec:eval-live})?

\subsection{Protocol and dataset}
\label{sec:eval-protocol}

\textbf{Ground truth.} The curated set contains 30 rows spanning
2022--2025, each with a primary public report, a verifiable routing
anomaly and (where the report allows it) a verifiable traffic anomaly:
15 outage rows (one of which is a negative control), 8 route leaks and
7 origin-hijack victims.

\textbf{What can be reproduced.} The 30 rows are scored from per-event
prediction files that ship with the artefact set, so every number in this
section can be recomputed row by row. What cannot be recomputed is the
evidence underneath them: the 26 cache bundles the rows cite are not in
the repository, and the 145 caches that are belong to a different event
corpus. The negative control and three further rows were scored on
synthetic timelines and are marked as such. We therefore report
\emph{archived} predictions throughout, name the file behind each table,
and provide \texttt{check\_eval\_provenance.py} (in \texttt{scripts/}),
which fails any artefact whose evidence reference does not resolve. A
held-out set of 9 further events is in the same position and is reported
separately in \S\ref{sec:eval-audit}.

\textbf{Strict protocol.} Three counting rules apply to every method in
every table. An abstention is never counted as correct; a method is scored
only on the events where it commits; and coverage (the share of events
where it commits) is reported next to accuracy, because a method that
commits twice and is right twice is not comparable to one that commits 24
times. Confidence intervals are Wilson intervals~\cite{wilson-1927}; paired
comparisons use the McNemar exact test. All Phase-1 numbers are produced
with the language model disabled.

A table is only as good as the artefact behind it, so we state the
artefact for every table and what the evidence check returns on it
(Appendix~\ref{app:provenance}). The check is mechanical: a script in
\texttt{scripts/} resolves every evidence reference, fails a row that names
none, fails a cache whose own ASN disagrees with the row, and exits
nonzero. A table passes only if every reference resolves. The
system-level comparison (Table~\ref{tab:baseline}) and the paired test
below it follow the same convention; the caches on disk belong to a
different event corpus, and matching them to the curated rows by ASN and
date resolves exactly one event (Orange Spain, 2024-01-03). The
directional layer is therefore measured on planted mechanisms and on
coverage, both in \S\ref{sec:eval-direction}.
\textbf{A timestamp defect in the cached corpus.} In part of the cached
set every event falls outside the cache's own window, and the modal residual
is exactly $-8$\,h: the window is written as UTC while the event log
carries local time. The bucket grid is anchored to the window, so those
events would be silently dropped and the event would look unevaluable.
\S\ref{sec:eval-direction} measures the direction test on this corpus with
and without a whole-hour repair (\texttt{agents/cache\_resolution.py}); the
repair recovers 15 of the 39 affected windows, and the rest still fail
because the stored timeline contains no routing events at all.

\subsection{Phase 1: cause classification}
\label{sec:eval-cause}

\begin{table}[t]
\centering
\caption{Ground-truth composition and strict Top-1 accuracy of the
Phase-1 arbitration (30 rows, all RRC00-verifiable). Precision, recall
and $F_1$ are computed per class from
\texttt{results/eval/v19\_live\_honest\_per\_event.csv}; the negative
control is counted as its own row and as a failure when the system
commits a cause.}
\label{tab:dataset}
\footnotesize
\setlength{\tabcolsep}{4pt}
\begin{tabular}{l r r r r r}
\hline
\textbf{Class} & \textbf{$n$} & \textbf{Top-1} & \textbf{Prec.} & \textbf{Rec.} & \textbf{$F_1$} \\
\hline
Outage (incl.\ 1 NC)  & 15 & 13/15 = 86.7\% & 0.867 & 0.867 & 0.867 \\
Route leak            &  8 &  7/8  = 87.5\% & 0.875 & 0.875 & 0.875 \\
Origin hijack victim  &  7 &  6/7  = 85.7\% & 1.000 & 0.857 & 0.923 \\
\hline
\textbf{Total}                & 30 & \textbf{26/30 = 86.7\%} & --- & --- & \textbf{0.888}\\
\textbf{Total (excl.\ NC)}    & 29 & \textbf{26/29 = 89.7\%} & --- & --- & --- \\
\hline
held-out (9)     & +9 & 6/9 = 66.7\% & --- & --- & ---   \\
\hline
\end{tabular}
\end{table}

The miss set is one error in each cause class plus the negative control;
\S\ref{sec:eval-errors} analyses them individually.

\begin{table}[t]
\centering
\caption{System-level comparison on the 29 attributable ground-truth rows. The
negative control is excluded here because it carries no cause label; it
is scored as a failure of the full system in Table~\ref{tab:dataset}.
Abstention is never counted as correct.}
\label{tab:baseline}
\small
\begin{tabular}{l r r r}
\hline
\textbf{System} & \textbf{Top-1} & \textbf{F1} & \textbf{Caus.\,F1} \\
\hline
Pick-Max (single detector)          & 34.5\% & 0.28 & 0.21 \\
Routing-only (no traffic signal)    & 51.7\% & 0.41 & 0.31 \\
Traffic-only (no routing signal)    & 41.2\% & 0.46 & 0.34 \\
\hline
\textbf{BGP Tracer (full)}          & \textbf{89.7\%} & \textbf{0.866} & \textbf{0.91} \\
\hline
\end{tabular}
\end{table}

The arbitration contributes 55.2\,pp over Pick-Max on Top-1 and 0.58 on
F1. We report no latency comparison against the single-signal variants:
those runs were timed on a separate 23-event run, so the columns would not
be comparable. A gradient-boosted baseline trained on the curated labels
is likewise absent, because at 29 events the training and evaluation folds
cannot be separated, and the resulting number would measure memorisation of
the label set rather than generalisation.

The contribution of each detector is visible in how it behaves across
classes rather than in a separate ablation run: the traffic detector stays
silent on origin-hijack victims, because for a hijack victim the traffic
drop is the hijack's effect rather than a co-incident outage, so its
silence is what separates the two classes.

Every comparison in this section is paired at the event level, so the
appropriate test is the McNemar exact two-sided test. The only significant
comparison on this set is the Phase-1 arbitration against Pick-Max
($p=0.003$): of the 29 events on which the two methods disagree in
outcome, BGP Tracer is correct on 21 and Pick-Max on 5. We report no
direction-level $p$-value, because a paired test over directions needs two
commitments per event and the surviving evidence does not supply them;
direction is reported through size, power and coverage in
\S\ref{sec:eval-direction} instead.

\subsection{Phase 2: operating characteristics of the direction test}
\label{sec:eval-direction}

A directional test has two properties worth measuring: whether it recovers
a mechanism that is actually there, and what it does on windows where the
answer is unknown. The first is measured with planted mechanisms, which can
be generated exactly and at scale; the second by running the same procedure
over every cache in the tree and counting commitments. We report both, size
and power from the first, coverage from the second.

\paragraph{Size.} The size experiment plants no effect---two
independent count series with no routing edge---and records how often the
single V2 test rejects; Appendix~\ref{app:aea-size} reports it over 400
windows per level.

The test is conservative at independence and anti-conservative under serial
dependence: at $\rho=0.8$ it rejects $12.3\%$ of the time at a nominal
$5\%$. The placebo null assumes the response is exchangeable across onsets
and a correlated series is not, so the placebo spread is too narrow. We
report the number rather than the nominal level, and we note the direction
of the error: it inflates commitments in exactly the
regime---volatile, correlated traffic---where a spurious direction is most
plausible. This is one reason the coverage of Table~\ref{tab:aea-coverage},
and not power, is the binding constraint on what the method can be used for.

\paragraph{Power.} Table~\ref{tab:aea-power} plants a routing pulse and a
traffic drop of the same length, separated by a known lag.

\begin{table}[t]
\centering
\caption{End-to-end power: share of 200 windows in which a direction is
committed, and the accuracy of those commitments. The planted mechanism
is a routing pulse followed by a traffic drop at the stated lag; nothing
else about the two series is known to the test.}
\label{tab:aea-power}
\small
\begin{tabular}{l r r r}
\hline
\textbf{Lag} & \textbf{Drop 20\%} & \textbf{Drop 50\%} & \textbf{Drop 80\%} \\
\hline
2  & 153/200 & 200/200 & 200/200 \\
4  & 184/200 & 200/200 & 200/200 \\
8  & 195/200 & 200/200 & 200/200 \\
16 & 159/200 & 200/200 & 200/200 \\
\hline
\multicolumn{4}{@{}l}{\footnotesize Accuracy among commitments: $100\%$ in all twelve cells.} \\
\end{tabular}
\end{table}

Power is high and, more important, every commitment is correct: across the
twelve cells the test committed $2{,}291$ times and never named the wrong
direction. The weakest cell is a small drop two buckets after the cause
($153/200$); a $20\%$ drop arriving sixteen buckets later is slightly
harder than one arriving at eight, because the post window of
Definition~\ref{def:segmentation} then covers only part of the pulse. The
non-monotonicity along the drop axis is sampling noise across $200$
windows, not a property of the rule.

\paragraph{Intervals versus point estimates.} The ordering rule is the
paper's central methodological claim, so we measure it directly against the
heuristic it replaces. Both orderings face the same remaining vetoes; only
the way the onset is reported differs. The point comparator reports a single
BIC-optimal split and commits whenever the two points differ; the shipped
rule reports the set of admissible split positions and commits only when
those sets are disjoint. The generator plants a \emph{simultaneous} pulse in
both series, where no direction exists, so every commitment is a false
commitment (\texttt{scripts/aea\_point\_vs\_interval.py}).

\begin{table}[t]
\centering
\caption{False commitments on a simultaneous common shock. Both rows
commit a direction only when the two onsets differ, so the difference is
the interval rule, not a different threshold. \emph{Point} reports a single
BIC-optimal split; \emph{interval} reports the set of admissible splits
and commits only when the two sets are disjoint.}
\label{tab:aea-point}
\small
\begin{tabular}{l r r r}
\hline
\textbf{Shock} & \textbf{Windows} & \textbf{Point} & \textbf{Interval} \\
\hline
Simultaneous pulse           & 600 & 9.7\% & 0.7\% \\
Simultaneous, 1-bucket skew  & 600 & 99.8\% & 94.7\% \\
\hline
\end{tabular}
\end{table}

Table~\ref{tab:aea-point} shows both halves of the result. On a truly
simultaneous shock the interval rule cuts false commitments by an order of
magnitude, from $9.7\%$ to $0.7\%$. On a shock skewed by a single bucket
neither rule can help: a one-bucket offset is a real ordering, so the
interval is wide enough to contain it and the difference between the two
orderings collapses. We report the second row because it is the boundary of
the claim: the interval rule protects against \emph{simultaneity}, not
against near-simultaneity, and it cannot do so without a resolution the data
do not support.

\paragraph{Coverage on the corpus.} Table~\ref{tab:aea-coverage} applies the
procedure to every cache in the tree.

\begin{table}[t]
\centering
\caption{Real-data coverage over the cached windows. ``Two series
built'' counts windows that yield two non-empty bucket series; the rows
above it are the reasons a window cannot be scored at all, which is where
the volume is. The final row reruns the measurement with $\alpha=0.99$, so
that V2 can never reject, to test whether abstention is a threshold
artefact.}
\label{tab:aea-coverage}
\footnotesize
\begin{tabular}{l r}
\hline
\textbf{Outcome} & \textbf{Windows} \\
\hline
Empty routing buckets        & 116 \\
No events inside the window  & 44 \\
Empty traffic buckets        & 37 \\
No declared window           &  3 \\
\addlinespace
Two series built             & 27 \\
\quad stopped at V1          & 27 \\
\quad committed              &  0 \\
\addlinespace
Committed at $\alpha=0.99$   &  0 \\
\hline
\end{tabular}
\end{table}

Two facts in that table locate the constraint. First, it sits upstream of
the test: most windows cannot produce both series at all, because the cached
aggregate holds no routing events in its buckets, so the test is never
reached. Second, the zero is not a threshold artefact: with $\alpha=0.99$,
which disables V2 entirely, the count is still zero, because every one of
the windows that produced two series stops at V1, twenty of them with no
routing edge, six with no traffic edge, and one with overlapping onsets. A
method that answered on these windows would be inventing the evidence, and
relaxing a level would not change the answer. Both runs are recorded under
protocol hash \texttt{541554a473d4} in
\texttt{results/eval/aea\_coverage\_caches.json}, together with the veto
tally.

\subsection{Input shape: the same window under two shapes}
\label{sec:eval-shape}

The contract of \S\ref{sec:impl-contract} is not an implementation detail,
and this section measures what it changes. We take one real cached
window---the Orange Spain RPKI-credential hijack of 2024-01-03
(\S\ref{sec:case})---and feed its evidence bundle to the engine twice:
once in the shape the live path produces, a merged event timeline, and once
with the same records shaped by Definition~\ref{def:grid}.

The two shapes differ in more than bucketing, and the difference is the
finding. The routing block holds $9{,}158$ announcement-level records, but
the merged timeline retains forty-nine of them at a single timestamp plus
one synthesized traffic record, so whichever grid it is placed on the
result has two non-filler points: the engine's five-minute bucketing yields
$91$ buckets, and a $30$-second grid---the finest the bucketing code
allows, capped at $120$ buckets---puts the routing records in the final
bucket and the traffic record in the first. In both cases the two points
are 7.5 hours apart, $89$ or $118$ buckets are empty and floored at
$\varepsilon=10^{-9}$, and neither series has an interior. A grid with two
non-filler points is not a time series; it is two spikes and a floor.

Under the contract the engine reads the second shape and refuses at the
first veto: with a single non-filler point on the routing side there is no
onset interval to order, so V1 returns
\texttt{inadmissible\_precedence} with sub-verdict
\texttt{onset\_unresolved} (\S\ref{sec:case} gives the statistics the
channels would have reported). For this window the refusal is also the
correct answer, because the ground truth records no traffic anomaly, but
the engine declines because the routing evidence collapsed into a single
bucket rather than because it observed an empty data plane. On a window
with the same shape and the opposite cause the two readings would not
coincide, which is why \S\ref{sec:method-grid} treats the input shape as
part of the claim. The evaluated windows of \S\ref{sec:eval-direction} are
not necessarily of this shape: the replay path reads each cache's full
event streams where they exist and falls back to the merged timeline only
where they do not. The contract is pinned by regression tests that fail if
either shape is misread: a raw event log must yield a bucketed series with
quiet gaps rather than a near-constant one, and padding a short window
must never manufacture buckets that read as ``measured, no anomaly''.

\paragraph{Decision latency.}
\label{sec:eval-latency}

We measure end-to-end latency on the alert replay rather than a per-stage
decomposition, because the first-stage screen in the evaluated build is a
fixed-threshold rule and a per-stage number attributed to it would
characterise that rule rather than the staged architecture. Over the $130$
completed replays the median end-to-end time is $225.6$\,s, the 95th
percentile is $783.7$\,s and the maximum is $1039.0$\,s
(Table~\ref{tab:latency} in Appendix~\ref{app:live-replay}). This is slower
than a one-minute triage target and we do not claim otherwise. The tail is
dominated by the Cloudflare Radar fetch and by snapshot loads for the path
and AS-relationship models, so shortening it requires a different fetching
strategy, not a different estimator.

\subsection{Error analysis}
\label{sec:eval-errors}

The Phase-1 miss set contains three errors, one per cause class, plus the
negative control.

\emph{E17, a route leak classified as an origin hijack.} The traffic
detector fired on a coincident content shift, and the arbitration read a
traffic anomaly with a leak-like path as a hijack. A traffic-side check
that separates a coincident content shift from a leak effect would be
needed to fix it, and the Phase-2 engine cannot supply one: the engine
abstains on this event's real cache, for the reasons in
\S\ref{sec:eval-direction}.

\emph{E23, a small-magnitude outage classified as normal.} The event fell
below the regularity threshold of the traffic detector. The routing-only
fallback (\S\ref{sec:method-phase1}) is the intended mitigation: a
withdrawal regime is a control-plane fact and does not need the traffic
baseline.

\emph{E31, a multi-party leak attributed to the wrong AS.} The cascade scope
was ambiguous between two candidate origins. This is a resolution limit of
the current AS-relationship snapshot rather than a rule error, and the
mitigation is cascade-scope expansion.

The negative control is a different kind of failure: the system committed a
cause where the correct answer is ``no clear anomaly''. We keep it as its own
row rather than folding it into a footnote because it is the only event
that tests the fourth verdict class, and because Phase~1 has no abstention
mechanism of its own---a property we return to in \S\ref{sec:discussion}.

\paragraph{Held-out set, narrative audit and channel sensitivity.}
\label{sec:eval-audit}

\textbf{The language model is a renderer.} Phase~2 renders reports with a
language model, but the model never overwrites a Phase-1 verdict. A factual
audit on nine held-out events (Table~\ref{tab:llm-audit}) shows high ASN
recall and perfect verdict consistency, and near-zero recall for numeric
traffic percentages.

\begin{table}[t]
\centering
\caption{Language-model narrative factual audit on 9 held-out events.}
\label{tab:llm-audit}
\small
\begin{tabular}{l r r}
\hline
\textbf{Dimension} & \textbf{Mean recall} & \textbf{Median recall} \\
\hline
ASN                 & 83.3\% & 100\% \\
Date                & 44.4\% & 0\%   \\
Traffic pct         & 0.0\%  & 0\%   \\
Integer intervals   & 22.2\% & 25\%  \\
\hline
Verdict consistency & \multicolumn{2}{c}{\textbf{100\% (3/3)}} \\
\hline
\end{tabular}
\end{table}

We report this because it is the empirical justification for the paper's
division of labour: the model may describe a verdict, and it may not
produce one. A model with 0\,\% recall on the quantity operators quote
first is a presentation risk, not an analyst.

\textbf{Held-out set.} The nine held-out events were collected after the
design was frozen and are independent of the design labels. The system
attains 6/9 = 66.7\,\% Top-1, a 23-point drop from the curated set. Three
structural factors account for the gap: distribution shift (the held-out
events favour longer-tail causes the curated set does not represent), sparse
traffic coverage (five of the nine land on ASes without Cloudflare Radar
coverage, so the strongest discriminator is absent), and stale routing caches
(two events have a routing cache that predates the event by more than 30
days). Each is an identifiable and addressable failure mode rather than a
property of the method, and each maps to a concrete data requirement:
broader radar coverage, fresher snapshot pinning, and a curated set
stratified by cause tail.

\textbf{Channel sensitivity.} Removing MSBD, CTPR, CIC, Granger or transfer
entropy individually changes the committed direction on \textbf{0\,\%} of
events, whereas removing CSS changes it on 37.9\,\%. Neither number is a
robustness statement about the channels. Zero flips means those five
channels never vote independently: they read the same two series, and before
the audit MSBD, CTPR and CIC also consumed the \texttt{event\_type} label,
so they were five views of one signal rather than five pieces of evidence.
The 37.9\,\% flip rate measures how much the earlier CSS-anchored prior
contributed, not how much causal information the CSS statistic carries. Both
observations motivated the uniform prior, the clustering and the
geometric-mean pooling of \S\ref{sec:method-gate}, and they are the reason a
channel-level ablation is reported here as a dependence diagnostic rather
than as a robustness claim.

\subsection{Live deployment}
\label{sec:eval-live}

BGP Tracer ran on a public alert pipeline for five months, consuming RIPE
RIS RRC00 updates and Cloudflare Radar five-minute aggregates on the same
orchestration as the ground-truth evaluation. The operational value of a
five-month run is throughput durability: the arbitration sustains production
rates for months without operator intervention, and the cache the live
pipeline writes is the cache the offline evaluation reads.

The replay denominator is the set of alerts that the pipeline actually
adjudicated. Of the 192 alerts replayed, 130 reached a verdict, 14 were
adjudicated \texttt{real\_incident} and 116 were not, giving a precision of
10.8\,\% (Wilson 95\,\% CI 6.5--17.3\,\%) and a suppression ratio of 9.3:1
on the processed stream. We report precision on that stream rather than a
pooled suppression ratio over all 192, because the remaining alerts did not
receive a verdict and counting them would mix collection failures into
arbitration quality. For the same reason the ratio is not compared with
operator-facing alarm counts: what the detectors emit are per-event (prefix,
timestamp) observations, and an operator alarm is a different unit.

The 14 confirmed incidents split evenly across the two sources whose ground
truth is genuinely open: 7 came from Cloudflare's curated annotation feed and
7 from the live hijack-monitoring feed. The latter include the AS1197
cross-ISP route leak of 2025-06-11, which affected 141 victim ASes across 28
countries, and the AS834 hijack of 2025-06-25, which touched victims in
eight countries. The split matters because it shows that the system fires on
incidents it discovers in a hijack feed, not only on incidents that the
curated feed already annotates.
